% GENERATED FILE. Edit canonical lessons, then run npm run generate:books.
% canonical-source-hash: fnv1a64:65e9c145c324a21d
% canonical-lessons: IT-C336-preriscaldare, IT-C336-infornare, IT-C336-sfornare, IT-C336-teglia, IT-C336-cottura, IT-R336-ripasso-329-336

\chapter{In forno}
\label{ch:it-in-forno}
\hlchaptermodality{336}

\begin{chapteropening}
\textbf{By the end of this chapter:} \emph{I can describe baking: preheating, putting in and taking out, the baking tray and the cooking time.}

Complete the last lesson of chapter 336: I can describe baking: preheating, putting in and taking out, the baking tray and the cooking time.
\end{chapteropening}

\section[preriscaldare]{preriscaldare — to preheat}

\label{lesson:IT-C336-preriscaldare}



\begin{quote}
Before the new one: say the Italian for to rise, then the Italian for the rolling pin.
\end{quote}

\subsection*{You'll want to know: preriscaldare}
\textbf{preriscaldare} — \textquotedblleft{}to preheat\textquotedblright{}.

\begin{quote}
Preriscalda il forno a 180 gradi.
\end{quote}

\begin{quote}
\emph{Preheat the oven to 180 degrees.}
\end{quote}

\emph{Il forno preriscaldato} — the preheated oven.

\begin{cousinweb}
\emph{preriscaldare} joins \emph{pre-} \textquotedblleft{}before\textquotedblright{}, \emph{ri-} and \emph{scaldare} \textquotedblleft{}to warm\textquotedblright{}, from Late Latin \emph{excaldāre}, built on \emph{calidus} \textquotedblleft{}hot\textquotedblright{} (Lewis \& Short).
\end{cousinweb}

\begin{grammarlens}[title={What you've built}]
One more word for telling, explaining and describing at B1.
\end{grammarlens}

\subsection*{Your turn}
\begin{itemize}
\raggedright
  \item \emph{Say it:} \emph{preriscaldare}
  \item \emph{Say it:} \emph{preriscaldare}, once more
  \item \emph{Say it:} a sentence of your own with \emph{preriscaldare}
  \item \emph{Recall:} say the Italian for to rise, then the Italian for the rolling pin, then say \emph{preriscaldare} again
\end{itemize}

\begin{tcolorbox}[breakable,colback=teal!4,colframe=teal!35!black,title={Before you move on}]
What does preriscaldare mean? (\textquotedblleft{}To preheat\textquotedblright{}.) Say it once more.
\end{tcolorbox}
\section[infornare]{infornare — to put in the oven}

\label{lesson:IT-C336-infornare}



\begin{quote}
Before the new one: say the Italian for the rolling pin, then the Italian for to preheat.
\end{quote}

\subsection*{You'll want to know: infornare}
\textbf{infornare} — \textquotedblleft{}to put in the oven\textquotedblright{}.

\begin{quote}
Inforna la torta per quaranta minuti.
\end{quote}

\begin{quote}
\emph{Bake the cake for forty minutes.}
\end{quote}

Built on \emph{forno}, ‘oven’.

\begin{cousinweb}
\emph{infornare} is built on \emph{forno} \textquotedblleft{}oven\textquotedblright{}, inherited from Latin \emph{furnus} (Lewis \& Short).
\end{cousinweb}

\begin{grammarlens}[title={What you've built}]
One more word for telling, explaining and describing at B1.
\end{grammarlens}

\subsection*{Your turn}
\begin{itemize}
\raggedright
  \item \emph{Say it:} \emph{infornare}
  \item \emph{Say it:} \emph{infornare}, once more
  \item \emph{Say it:} a sentence of your own with \emph{infornare}
  \item \emph{Recall:} say the Italian for the rolling pin, then the Italian for to preheat, then say \emph{infornare} again
\end{itemize}

\begin{tcolorbox}[breakable,colback=teal!4,colframe=teal!35!black,title={Before you move on}]
What does infornare mean? (\textquotedblleft{}To put in the oven\textquotedblright{}.) Say it once more.
\end{tcolorbox}
\section[sfornare]{sfornare — to take out of the oven; to turn out (lots of something)}

\label{lesson:IT-C336-sfornare}



\begin{quote}
Before the new one: say the Italian for to preheat, then the Italian for to put in the oven.
\end{quote}

\subsection*{You'll want to know: sfornare}
\textbf{sfornare} — \textquotedblleft{}to take out of the oven; to turn out (lots of something)\textquotedblright{}.

\begin{quote}
Ho appena sfornato il pane.
\end{quote}

\begin{quote}
\emph{I have taken the bread out of the oven this minute.}
\end{quote}

Also figurative: \emph{sforna un libro all'anno}.

\begin{cousinweb}
\emph{sfornare} is \emph{s-} \textquotedblleft{}out of\textquotedblright{} with \emph{forno} \textquotedblleft{}oven\textquotedblright{}, from Latin \emph{furnus} (Lewis \& Short).
\end{cousinweb}

\begin{grammarlens}[title={What you've built}]
One more word for telling, explaining and describing at B1.
\end{grammarlens}

\subsection*{Your turn}
\begin{itemize}
\raggedright
  \item \emph{Say it:} \emph{sfornare}
  \item \emph{Say it:} \emph{sfornare}, once more
  \item \emph{Say it:} a sentence of your own with \emph{sfornare}
  \item \emph{Recall:} say the Italian for to preheat, then the Italian for to put in the oven, then say \emph{sfornare} again
\end{itemize}

\begin{tcolorbox}[breakable,colback=teal!4,colframe=teal!35!black,title={Before you move on}]
What does sfornare mean? (\textquotedblleft{}To take out of the oven; to turn out (lots of something)\textquotedblright{}.) Say it once more.
\end{tcolorbox}
\section[la teglia]{la teglia — the baking tray, the baking tin}

\label{lesson:IT-C336-teglia}



\begin{quote}
Before the new one: say the Italian for to put in the oven, then the Italian for to take out of the oven.
\end{quote}

\subsection*{You'll want to know: la teglia}
\textbf{la teglia} — \textquotedblleft{}the baking tray, the baking tin\textquotedblright{}.

\begin{quote}
Metti le patate nella teglia.
\end{quote}

\begin{quote}
\emph{Put the potatoes in the baking tray.}
\end{quote}

\emph{La pizza in teglia} — pizza baked in a tray.

\begin{cousinweb}
\emph{teglia} is inherited from Latin \emph{tēgula} \textquotedblleft{}a tile\textquotedblright{} (EtymDB), from \emph{tegere} \textquotedblleft{}to cover\textquotedblright{}.
\end{cousinweb}

\begin{grammarlens}[title={What you've built}]
One more word for telling, explaining and describing at B1.
\end{grammarlens}

\subsection*{Your turn}
\begin{itemize}
\raggedright
  \item \emph{Say it:} \emph{la teglia}
  \item \emph{Say it:} \emph{la teglia}, once more
  \item \emph{Say it:} a sentence of your own with \emph{la teglia}
  \item \emph{Recall:} say the Italian for to put in the oven, then the Italian for to take out of the oven, then say \emph{la teglia} again
\end{itemize}

\begin{tcolorbox}[breakable,colback=teal!4,colframe=teal!35!black,title={Before you move on}]
What does la teglia mean? (\textquotedblleft{}The baking tray, the baking tin\textquotedblright{}.) Say it once more.
\end{tcolorbox}
\section[la cottura]{la cottura — the cooking; the cooking time}

\label{lesson:IT-C336-cottura}



\begin{quote}
Before the new one: say the Italian for to take out of the oven, then the Italian for the baking tray.
\end{quote}

\subsection*{You'll want to know: la cottura}
\textbf{la cottura} — \textquotedblleft{}the cooking; the cooking time\textquotedblright{}.

\begin{quote}
La cottura dura un'ora.
\end{quote}

\begin{quote}
\emph{The cooking takes an hour.}
\end{quote}

\emph{Il tempo di cottura} — cooking time; \emph{a metà cottura} — halfway through.

\begin{cousinweb}
\emph{cottura} is inherited from Latin \emph{coctūra} \textquotedblleft{}a cooking\textquotedblright{} (EtymDB), from \emph{coquere} \textquotedblleft{}to cook\textquotedblright{}.
\end{cousinweb}

\begin{grammarlens}[title={What you've built}]
One more word for telling, explaining and describing at B1.
\end{grammarlens}

\subsection*{Your turn}
\begin{itemize}
\raggedright
  \item \emph{Say it:} \emph{la cottura}
  \item \emph{Say it:} \emph{la cottura}, once more
  \item \emph{Say it:} a sentence of your own with \emph{la cottura}
  \item \emph{Recall:} say the Italian for to take out of the oven, then the Italian for the baking tray, then say \emph{la cottura} again
\end{itemize}

\begin{tcolorbox}[breakable,colback=teal!4,colframe=teal!35!black,title={Before you move on}]
What does la cottura mean? (\textquotedblleft{}The cooking; the cooking time\textquotedblright{}.) Say it once more.
\end{tcolorbox}
\section[(dialogue)]{Ripasso dei capitoli 329–336}

\label{lesson:IT-R336-ripasso-329-336}



\begin{quote}
Say the Italian for the revolt and to rebel.
\end{quote}

\subsection*{Your turn}
Say these aloud:

\begin{itemize}
\raggedright
  \item the revolt, to rebel, to invade, to surrender, to free
  \item the age, the Middle Ages, the ruins, the excavations, the archaeologist
  \item the ingredients, the amount, the pinch, the teaspoon, the bowl
  \item the chopping board, to slice, to chop finely, to peel, to grate
  \item to boil, to cook in boiling water, to drain, the lid, the shallow pan
  \item to fry gently, to brown, to fry, to deglaze, the soffritto
  \item to knead, the dough, the yeast, to rise, the rolling pin
  \item to preheat, to put in the oven, to take out of the oven, the baking tray, the cooking
\end{itemize}

\begin{tcolorbox}[breakable,colback=teal!4,colframe=teal!35!black,title={Before you move on}]
The revolt? (\textbf{la rivolta}.) To rebel? (\textbf{ribellarsi}.) To invade? (\textbf{invadere}.) To surrender? (\textbf{arrendersi}.) To free? (\textbf{liberare}.) The age? (\textbf{l'epoca}.) The Middle Ages? (\textbf{il medioevo}.) The ruins? (\textbf{le rovine}.) The excavations? (\textbf{gli scavi}.) The archaeologist? (\textbf{l'archeologo}.) The ingredients? (\textbf{gli ingredienti}.) The amount? (\textbf{la dose}.) The pinch? (\textbf{il pizzico}.) The teaspoon? (\textbf{il cucchiaino}.) The bowl? (\textbf{la ciotola}.) The chopping board? (\textbf{il tagliere}.) To slice? (\textbf{affettare}.) To chop finely? (\textbf{tritare}.) To peel? (\textbf{sbucciare}.) To grate? (\textbf{grattugiare}.) To boil? (\textbf{bollire}.) To cook in boiling water? (\textbf{lessare}.) To drain? (\textbf{scolare}.) The lid? (\textbf{il coperchio}.) The shallow pan? (\textbf{il tegame}.) To fry gently? (\textbf{soffriggere}.) To brown? (\textbf{rosolare}.) To fry? (\textbf{friggere}.) To deglaze? (\textbf{sfumare}.) The soffritto? (\textbf{il soffritto}.) To knead? (\textbf{impastare}.) The dough? (\textbf{l'impasto}.) The yeast? (\textbf{il lievito}.) To rise? (\textbf{lievitare}.) The rolling pin? (\textbf{il mattarello}.) To preheat? (\textbf{preriscaldare}.) To put in the oven? (\textbf{infornare}.) To take out of the oven? (\textbf{sfornare}.) The baking tray? (\textbf{la teglia}.) The cooking? (\textbf{la cottura}.)
\end{tcolorbox}
